\section{Testing the Native Operator Link}
\label{v480:native-operator-link}

The previous interface becomes predictive only if the native projector,
pacer and generation embedding determine its matrix elements. The
retained definitions provide a useful partial link, but not that full
identification. All statements below use the same positive-metric,
bounded or finite realization as Section~\ref{v480:generation-coupling}.
Orthogonality of $P$, common domains and the seed metric are assumptions
of this realization, not consequences of the inherited idempotence
condition alone. The HC projector $C$ remains distinct from $P$ and
from the defect penalty $\mathcal H_C$.

\subsection{What the Original Definitions Supply}
\label{v480:native-source-link}

The Meta-Septum definition before Theorem~\ref{thm:dirac-probe}
specifies an admissible projector through a trace condition modulo HC;
it does not list its entries on the generation seeds. That theorem
writes $R_\beta=e^{i\pi\bsym}R_{\mathrm{GUE}}$, without specifying
the clock eigenvectors relative to $P$ and those seeds.
Proposition~\ref{prop:equivalent-closure-criteria} requires
$[\mathcal M_{16},P]=0$ for exact Millennium-matrix closure.
No equation there identifies $R_{\mathrm{GUE}}$ with
$\mathcal M_{16}$ or a function of it. The HS(4)--HS(6) carrier
coding and the spectral-to-mass interface of
Proposition~\ref{prop:spectral-to-mass-realization-map} do not
evaluate these joint matrix elements either. A coordinate coding
does not specify an operator action.

\begin{proposition}[The two minimal penalty terms coincide]
\label{v480:native-penalty-identity}
Write $M=\mathcal M_{16}$ on the declared common carrier and
$P^\perp=I-P$. Set
\[
 K=[M,P],\qquad E=P^\perp MP+PMP^\perp,
 \qquad H=\rho K^*K+\sigma E^*E,
 \quad \rho,\sigma\geq0.
\]
Then
\[
 E^*E=K^*K,\qquad H=(\rho+\sigma)K^*K,
 \qquad [H,P]=0.
\]
For an isometric closed embedding $S$, let
$B_M=P^\perp MS/r_{\rm ref}$ and $Q_M=B_M^*B_M$, where
$r_{\rm ref}>0$ has the units of $M$. Then
\[
 S^*HS=(\rho+\sigma)r_{\rm ref}^2Q_M.
\]
\end{proposition}
\begin{proof}
In the orthogonal splitting $\operatorname{im}P\oplus\ker P$,
write $M=\left(\begin{smallmatrix}A&F\\G&D\end{smallmatrix}\right)$.
Then
\[
 K=\begin{pmatrix}0&-F\\G&0\end{pmatrix},\quad
 E=\begin{pmatrix}0&F\\G&0\end{pmatrix},\quad
 K^*K=E^*E=\begin{pmatrix}G^*G&0\\0&F^*F\end{pmatrix}.
\]
This proves the first claims without requiring $M=M^*$.
Since $PS=S$, $KS=P^\perp MS=r_{\rm ref}B_M$;
compressing the positive product proves the last identity.
\end{proof}

Thus the existing minimal defect penalty already has a Gram-type
compression. Identifying it with the pacer Gram matrix additionally
requires, for example, $R_\beta=e^{i\varphi}M$ with matched units.
That bridge is not supplied by the source. Even when $H\ne0$,
$P^\perp HP=0$: this minimal penalty cannot itself be the missing
off-sector pacer channel. The statement does not concern a separately
postulated deeper HC generator.

\begin{proposition}[Transferring closure to the clock requires a bridge]
\label{v480:native-closure-transfer}
If $[M,P]=0$ and $R_0$ belongs to the norm-closed unital
$*$-algebra generated by $M$, then $[R_0,P]=0$ and
$Q_\beta=0$ for $R_\beta=e^{i\varphi(\beta)}R_0$.
Exact closure of $M$ alone does not imply $Q_\beta=0$
for an independently specified pacer.
\end{proposition}
\begin{proof}
Taking adjoints gives $[M^*,P]=0$. Every finite product and sum
of $M,M^*,I$ therefore commutes with $P$, and so does every norm
limit. Its closed-to-complementary block vanishes. For the converse
warning, take $P=\operatorname{diag}(I_3,0)$,
$S=(I_3,0)^{\mathsf T}$, $M=I_6$ and
$R_0=\left(\begin{smallmatrix}0&I_3\\I_3&0\end{smallmatrix}\right)$.
Then $[M,P]=0$ but $Q_0=I_3$.
\end{proof}

In particular, identifying the clock as a polynomial or continuous
functional-calculus expression in an exactly closed bounded normal
$M$ kills this defect route. This does not prohibit nonzero masses
from the internal block $P\Omega P$ or another readout. It prohibits
deriving them from this particular transition while maintaining that
bridge and exact closure.

\subsection{Spectral and Frame Data Are Not Yet Joint Matrix Elements}
\label{v480:native-spectral-data}

\begin{proposition}[An unchanged clock spectrum can change the Gram spectrum]
\label{v480:native-isospectral-freedom}
Keep $P=\operatorname{diag}(I_3,0)$ and
$S=(I_3,0)^{\mathsf T}$. For real diagonal matrices
$D_+=\operatorname{diag}(d_i^+)$ and $D_-=\operatorname{diag}(d_i^-)$,
put
\[
 O_\theta=\begin{pmatrix}\cos\Theta&-\sin\Theta\\
                         \sin\Theta&\cos\Theta\end{pmatrix},
 \quad \Theta=\operatorname{diag}(\theta_i),\qquad
 R_\theta=O_\theta\operatorname{diag}(D_+,D_-)O_\theta^*.
\]
These self-adjoint clocks all have the same eigenvalues, whereas
\[
 Q_\theta=\frac{1}{4r_{\rm ref}^2}
 \operatorname{diag}\bigl((d_i^+-d_i^-)^2\sin^2(2\theta_i)\bigr).
\]
Thus even an exact simple clock spectrum does not by itself determine
the transition spectrum relative to fixed closure and seed spaces.
\end{proposition}
\begin{proof}
$O_\theta$ is orthogonal, so conjugation preserves the clock spectrum.
The lower-left block of $R_\theta$ is
$(D_+-D_-)\sin\Theta\cos\Theta$; its normalized Gram product gives
the formula. For nonzero gaps, varying the angles changes its
eigenvalues, not just its eigenvectors.
\end{proof}

This finite example is not a replacement for the native Riemann clock.
It identifies what a zero sequence or spacing statistic leaves
unspecified: the orientation of the clock eigenspaces relative to
the closure projector and generation embedding. A native joint
operator law could remove this freedom; its evaluation is the burden.

\begin{proposition}[A simultaneous unitary frame change preserves mass ratios]
\label{v480:native-frame-covariance}
Let $F$ be unitary on the carrier, $V$ unitary on $\mathbb C^3$,
and $s>0$. Starting with $G_0=T_0^*P_0T_0>0$, set
\[
 P_f=FP_0F^*,\quad T_f=FT_0V,\quad
 \widetilde R_f=s e^{i\varphi}F\widetilde R_0F^*.
\]
Then the normalized embeddings and transition matrices obey
\[
 S_f=FS_0V,\qquad B_f=s e^{i\varphi}FB_0V,
 \qquad Q_f=s^2V^*Q_0V.
\]
Consequently nonzero spectral mass ratios and $K_{\rm spec}$ are
unchanged under the candidate mass law.
\end{proposition}
\begin{proof}
$G_f=V^*G_0V$, and its positive inverse root is
$V^*G_0^{-1/2}V$. Substitute into $S_f=P_fT_fG_f^{-1/2}$,
then into $B_f=(I-P_f)\widetilde R_fS_f$. The Gram identity follows.
Common positive scaling and unitary conjugation preserve the ratios.
\end{proof}

This is a statement about a declared positive Hilbert metric, not an
identification of a Lorentz boost with $F$. A common scalar Planck
calibration or common frame phase cannot generate a hierarchy change
in this lane. A relative motion of the projector, clock or seeds,
or a nontrivial relative channel, requires its own native law; it is
not supplied by assigning the same frame label to all three.

\begin{proposition}[The transition requires sufficient output rank]
\label{v480:native-channel-rank}
For $B:\mathbb C^3\to\ker P$ with
$\dim\operatorname{im}B\leq d$, one has
$\operatorname{rank}(B^*B)=\operatorname{rank}B\leq d$.
Three strictly positive candidate masses require three independent
transition outputs. A single scalar complementary channel produces
at most one positive eigenmass in this Gram readout.
\end{proposition}
\begin{proof}
$\langle v,B^*Bv\rangle=\|Bv\|^2$ implies
$\ker(B^*B)=\ker B$. Rank-nullity on the three-dimensional input
proves the claim.
\end{proof}

The three entries $(\Psi_1,\Psi_2,\Psi_\perp)$ in
Proposition~\ref{prop:triadic-two-probe} label probe channels,
not the three generation coordinates. An operator-valued channel
can have dimension greater than one; the proposition does not prove
that the native complementary space is one-dimensional. An explicit
embedding and channel rank are needed before this block system can
be used as a generation coupling.

\subsection{A Conditional Mass Law From the Existing Block Architecture}
\label{v480:native-block-mass-law}

The inherited field operator $\Omega$ and its triadic blocks suggest
testing elimination of a gapped complementary sector. The following
is finite spectral algebra at a frozen background, in units
$\hbar=c=1$. It additionally assumes a self-adjoint mass operator,
canonical kinetic normalization and a closed truncation to the chosen
three-dimensional seed space plus its retained complement. The source
does not establish those assumptions or the required block entries.
Other closed states, derivative couplings and background variations
cannot silently be omitted. Hermiticity requires adjoint relations
between the paired triadic channels, not merely their names.

\begin{proposition}[Elimination of a gapped complementary block]
\label{v480:native-schur-mass}
Let the supplied mass-dimension operator on that truncation be
\[
 \Omega_{\rm tr}=\begin{pmatrix}A&V^*\\V&H_\perp\end{pmatrix},
 \qquad A=A^*,\quad H_\perp=H_\perp^*\geq\Delta I,
 \quad \Delta>0.
\]
For an eigenvalue $\lambda$ with $|\lambda|<\Delta$, its closed
coordinate $x$ satisfies
\[
 \bigl(A-V^*(H_\perp-\lambda I)^{-1}V\bigr)x=\lambda x,
 \qquad y=-(H_\perp-\lambda I)^{-1}Vx.
\]
Writing $A_{\rm eff}(\lambda)$ for the operator in parentheses,
\[
 \|A_{\rm eff}(\lambda)-A_{\rm eff}(0)\|
 \leq\frac{\|V\|^2|\lambda|}{\Delta(\Delta-|\lambda|)}.
\]
\end{proposition}
\begin{proof}
Solve the lower block equation for $y$ and substitute into the
upper one. The resolvent identity gives
$(H_\perp-\lambda I)^{-1}-H_\perp^{-1}
=\lambda(H_\perp-\lambda I)^{-1}H_\perp^{-1}$.
The spectral gap bounds the two inverse norms by
$(\Delta-|\lambda|)^{-1}$ and $\Delta^{-1}$.
\end{proof}

\begin{proposition}[When a Gram mass law is a leading approximation]
\label{v480:native-flat-gap-mass}
Assume further $A=0$, $H_\perp=\Delta I$ and $V=\mu B$, with
$\mu>0$ a mass scale, $B$ dimensionless and $Q=B^*B$.
For each $q_i\geq0$ the light eigenvalue is $-m_i$, where
\[
 m_i=\frac{\sqrt{\Delta^2+4\mu^2q_i}-\Delta}{2}
     =\frac{\mu^2}{\Delta}q_i
       +O\!\left(\frac{\mu^4q_i^2}{\Delta^3}\right).
\]
The physical interpretation as a positive mass magnitude is conditional
on this canonical Hermitian Dirac-type reading. The signed effective
operator is negative, not $+(\mu^2/\Delta)Q$. For a normalized light
eigenstate its complementary weight is
\[
 \|y\|^2=\frac{m_i}{\Delta+2m_i}.
\]
\end{proposition}
\begin{proof}
In an eigenvector of $Q$, elimination gives
$\lambda=-\mu^2q_i/(\Delta-\lambda)$, so the light root of
$\lambda^2-\Delta\lambda-\mu^2q_i=0$ is the displayed negative
or zero root. Expansion for $\mu^2q_i/\Delta^2\ll1$ gives the approximation.
Moreover $\|y\|^2/\|x\|^2=\mu^2q_i/(\Delta+m_i)^2
=m_i/(\Delta+m_i)$; normalization gives the weight formula.
\end{proof}

This supplies a possible justification for $M=\kappa Q$ as a
\emph{leading mass-magnitude law}, with $\kappa=\mu^2/\Delta$,
if the native field blocks independently establish a flat positive
gap, zero direct block and $V=\mu B$. A non-flat gap gives a
weighted Gram matrix $V^*H_\perp^{-1}V$, not generally a scalar
multiple of $Q$; a nonzero $A$ also changes the spectrum. Even with
a flat gap, the exact eigenmasses must be tested anew for Koide.
No native values of $\mu,\Delta$ or the triadic matrix elements
have been obtained by this calculation.

There is also an exact-closure limitation: an eigenstate constrained
to the original $P$-sector has $y=0$, hence $Vx=0$. Nonzero masses
from the special case $A=0$ above therefore require dressed states
with nonzero complementary weight. They are not exactly $P$-closed.
For $\mu\|B\|/\Delta\ll1$ the weight is controlled, or one may define
a distinct spectral projector onto the light eigenspace. Neither
choice is the original strict closure constraint. Closure of $M$
alone, invariance under the full $\Omega$, and approximate closure
of a dressed state are different conditions.

\paragraph{Result of the native search.}
The existing penalty yields an exact algebraic Gram compression, and
the existing block architecture admits a conditional gapped-sector
mass mechanism. Neither evaluates the native tuple $(P,R_\beta,T)$.
The next model-specific datum is the cross block
$(I-P)R_\beta S$, together with a declared relation, if any, to
$(I-P)\Omega S$. A clock spectrum, common frame shift, scalar
length calibration or finite coding alone cannot provide it. The
three-generation rank, relative-channel law, physical gap, Koide
radius and absolute mass scale remain independent obligations.
